\begin{resumo}
 Em métodos ágeis, os critérios de aceitação (CA) definem quando uma história de usuário está pronta e, quando escritos em \textit{Gherkin}, podem ser verificados de forma automatizada. Escrevê-los à mão, porém, exige tempo e conhecimento do domínio, e os modelos de linguagem de grande porte (LLMs) passaram a ser usados como apoio a essa tarefa. Ainda não está claro, no entanto, se fornecer ao modelo o contexto do próprio projeto melhora os critérios gerados. Este trabalho tem como objetivo avaliar a acurácia de cenários em \textit{Gherkin} gerados por LLMs a partir de histórias de usuário, comparando a geração com e sem geração aumentada por recuperação (RAG). Na primeira etapa, foi conduzida uma revisão estruturada da literatura, que reuniu 25 estudos. A revisão mostrou que o tema é recente, com 80\% dos estudos publicados a partir de 2025, que a geração se apoia principalmente na escrita do \textit{prompt} e que a avaliação dos critérios gerados ainda é frágil: poucos estudos os comparam com um gabarito, e apenas dois adotam de fato uma estratégia empírica. Com base nessas lacunas, foi planejado um estudo experimental com um único fator, o uso de RAG, aplicado em desenho pareado a histórias de usuário dos projetos MeasureSoftGram e CAPJu. A acurácia será medida pelas métricas Acc@Hit e Acc@Cor, a partir dos objetivos de teste atômicos do gabarito e do julgamento de um comitê de três LLMs, e a diferença entre os tratamentos será verificada por um teste pareado. A execução do estudo e a análise dos resultados estão previstas para a segunda etapa do trabalho.

 \vspace{\onelineskip}

 \noindent
 \textbf{Palavras-chave}: critérios de aceitação; \textit{Gherkin}; modelos de linguagem de grande porte; geração aumentada por recuperação; \textit{Behavior-Driven Development}.
\end{resumo}
